\section{Civilizations of Space and Civilizations of Time}

The remaining traces, often still only in stone, of some of the great civilizations of the origins often contain a rarely noticed meaning. In the face of what remains of the most ancient Greco-Roman world and then that of Egypt, Persia, China, eventually reaching the mysterious and mute megalithic monuments scattered among deserts, lands, and forests as the final emergent and immobile vestiges of submerged and hidden worlds---and, as the boundary in the direction opposite to the course of history, including some forms of Medieval Europe: in the face of everything that comes from asking if the miraculous resistance to time of such civilizations, due to a favorable concourse of exterior circumstances, does not also include a symbolic meaning.

This feeling is reinforced if we focus primarily on the general character of the life of the civilizations to which in a great part such traces belong, that is on the general character of life called ``traditional''. It is a life that has maintained the same through centuries and generations, in an essential fidelity to the same principles, to the same type of institutions, to the same vision of the world; in the face of calamitous events, also susceptible to adaption and external modifications, but unalterable in its nucleus, in its animating principle, in its spirit, in its unanimity.

Such a world seems to bring us back especially to the Orient. If you think of what China and India were up until relatively recent times, and even the Japanese up to recently. But, in general, the more you recede back in time, the more you feel that living, universal, and powerful type of civilization, so that the Orient must be considered only as the part of the world where, through fortuitous circumstances, it was able to subsist longer and develop better than elsewhere. In such a type of civilization the law of time seems to be, in part, suspended. More than in time, these civilizations seem to have had life in \emph{space}. They had a ``timeless'' character.

These, according to the formulation current today, would therefore be the ``stable'', ``static'' or ultraconservative civilizations. In reality, they are the civilizations whose material vestiges even seem destined to live longer than all the monuments or all the ideal creations of the modern world that appear powerless to endure more than a half century and within which the words ``progressive'' or ``dynamism'' mean only a compliance to contingency, to the movement of an incessant change, of a rapid ascent and a likewise rapid, vertiginous decline. Processes that do not obey a true internal organic law, that are not contained within boundaries, that atomize themselves and take the hand of those from whom they were favoured: that is the characteristic of this different world, in all its sectors. That does not prevent it being made a criterion of the measure for everything which in the eminent sense should be called ``civilization'', in the field of a historiography for which the other types are judged to have arrogant and pejorative values.

In this regard, the misunderstanding is typical that confuses immobility with something that had a rather different meaning in traditional civilizations: unchangeability. These civilizations were civilizations of \emph{being}. Their strength was manifested precisely in their identity, in the victory they achieved over becoming, over ``history'', over change, over unformed flow. They are the civilizations that descended into the depths, beyond the fickle and treacherous waters, and established solid roots in those depths.

The opposition between modern and traditional civilizations can be expressed this way: \emph{modern civilizations are devourers of space, traditional civilizations were devourers of time}.

The former, or modern civilizations, are vertiginous from their frenzy of motion and spatial conquest, generators of an inexhaustible arsenal of mechanical means capable of reducing every distance, shortening every interval, contracting everything that is scattered in a multitude of places into a feeling of ubiquity. The excitement of a need to possess; dark anguish in the face of everything that is detached, isolated, deep, or distant; the impulse to self-expansion, to circulate, and to be associated and to find itself in every place, outside as well as inside. Science and technology, promoted by this irrational existential impulse, in their turn reinforce it, feed it, exacerbate it: trade, communications, ultrasonic speeds, radio, television, standardization, cosmopolitanism, internationalism, unlimited production, the American spirit, the ``modern'' spirit.

Quickly the matrix extends itself, reinforces itself, perfects itself. Terrestrial space has hardly any more mysteries. The pathways through the earth, water, sky are opened up. The human gaze has probed the most remote heavens, the infinitely great and the infinitesimally small. One no longer speaks of other lands but of other planets. On ours, action is carried like lightening far and wide. The tumultuous confusion of a thousand voices blend little by little into a uniform, atonal, impersonal rhythm. They are the latest effects of what was called Western ``Faustianism'', which is not unlike the same revolutionary myth in its various aspects including that technocratic formulation in the field of a decadent messianism.

In contrast, traditional civilizations were vertiginous precisely in their stability, their identity, their unstoppable and immutable existence in the midst of the current of time and history: so that they became capable of expressing a shadow of eternity even in sensible and tangible forms. They were islands in time, bastions in time; they acted as forces consuming time and history. Precisely because of this characteristic, it is imprecise to say that they ``were''--- it would be more correct to say that they simply \emph{are}. If they seem to draw back and almost vanish in the distance of a past which sometimes has even mythical traits, that is only the mirage to which anyone is necessarily subject, when he is transported by an irresistible current that increasingly distances him from the places of spiritual stability. Moreover, this idea corresponds exactly to the image of the ``double perspective'' given by an ancient traditional teaching: the ``immobile lands'' retreat and move for those who go with the waters, the waters move and retreat for those who reside solidly in the ``immobile lands''.

To understand this image, relating it not the physical plane but the spiritual plane, means to perceive also the right hierarchy of values, wherever our gaze is carried beyond the horizon which encloses our contemporaries. What seemed to be the past makes itself present, through its essential relation of historical forms (contingent as such) to metahistorical contents. What was qualified as ``static'' is seen to be saturated with enormous life. \emph{The others}, they are the fallen, the decentered. The philosophy of becoming, historicism, evolutionism, and so on appear as the intoxications of castaways, like the truth typical of those who flee (``\emph{where are you fleeing to, idiots}?'' Bernanos), and of those who do not have inner consistency and who do not know what inner consistency is and not even the source of every true height and every actual achievement---of the same achievement that cannot be limited to the intangible and the often invisible spiritual heights but were in fact expressed equally in epics, in cycles of civilizations that, precisely, even in their mute and dispersed vestiges of stone, seem to conceal something supertemporal and eternal. To which are also added traditional artistic, monolithic, rough, and powerful creations, alien from everything that is subjectivistic, often anonymous, almost extensions of the same elementary forces.

Finally we must recall what the conception of the time itself was in traditional civilizations: not an irreversible linear conception, but cyclic, in periods. From a combination of customs, rites and institutions characteristic both of higher civilizations and of echoes existing in some peoples described as ``primitive'' (in such regard, we recommend the same material gathered from the history of religions by Huber, Mauss, Eliade, and others) the firm intention is to bring time back to the origins (whence the cycle), in the sense of destroying what in it that is simply becoming, of restraining it, of making it express or reflect metahistorical, sacred, or metaphysical structures, often connected to myth. In such terms---as a ``moving image of eternity''---time acquired a value and a direction, not as ``history''. Returning to the origins meant a restoration, drawn to the fountain of perpetual youth, confirming spiritual stability, against temporality. The great cycles of nature evoked this attitude. ``Historical consciousness'', inseparable from the situation of ``modern'' civilizations, only sealed the fracture, the fall of man into temporality. But it is presented as an achievement of the last man, that is, of crepuscular man.

It is not unusual that, even within the field of presumed scientific objectivity, certain discoveries, the origins of general conceptions destined to revolutionize an epoch, have the character of a sign, so that their happening in a given period, and not in another, should not be considered accidental. Referring, for example, to the science of nature, it is more or less known to anyone, according to the latest voguish theories of Einstein and his disciples, that it is an indifferent thing to affirm that the Earth moves around the sun, or the opposite: it is only a question of preferring a greater of lesser complication of astrophysical calculations in the fixing the relational systems. Now, it seems rather significant that the Copernican discovery, with which it ceased to be ``true'' that the Earth is the solid and immobile center of the celestial entities, and it becomes ``true'' instead that it moves itself, and through its law is wandering in the cosmic space, as an irrelevant part of a dispersed system of expansion into the indefinite. It seems, let's say, rather significant that this discovery happened more or less at the time of the Renaissance and humanism, that is, at the time of the most decisive upheaval for the coming of the new civilization, in which the individual lost little by little every connection with that which ``is'', he degenerated from every spiritual centrality to the point of making his own the point of view of becoming, of history, of change, of the irrepressible and unpredictable current of life. The most remarkable upheaval was instead the pretense---the illusion---of having finally discovered, of affirming and glorifying ``man'', hence the term ``humanism''. In fact this was reduced to that which is ``only human'' with the impoverishment of the possibility of an opening and an integration into the ``more than human''.

This is not the only one of the symbolic revolutions that could be mentioned in this matter. Using the example we made, one point must be made precise about the ``Copernican revolution'': in the traditional world no so-called ``objective'' truth was important. Truths of that type could be also be incidentally considered, in the sense of their actual relativity on one hand, and of their human value on the other hand, bearing in mind the criteria of opportunities in regards to the general mode of feeling. A traditional theory of nature could therefore be also ``incorrect'' from the point of view of modern science (in one of its stages). But its value, the reason for which it was superior lay in its susceptibility to act as a means of expression of something true on a different and more interesting plane. For example, the geocentric theory was that the world of sensible appearances grasped an aspect suitable to act as a support for a truth of a different and unassailable type, to the truth regarding ``being'', the spiritual centrality, as principle of the true essence of man.

This will suffice to clarify morphologically the contrast between the civilizations of space and the civilizations of time. From that, it would also be easy to deduce the corresponding typological and existential contrast between the man of one civilization and the man of another. And if we are to move on to the problem of the crisis of the present epoch on the basis of what was already said, the uselessness of any criticism whatsoever, of any reaction, of any velleity of rectifying actions should appear in a rather clear way before an internal change of polarity, a metanoia, to use the ancient term, in man himself or, at least, in a determinate number of men in positions of exercising a decisive influence, is produced: in the sense of a shift toward the dimension of ``being'', of ``that which is'', a dimension gone missing and dissolved in modern man so that very little is any longer known about what internal stability and centrality are, and therefore also calm and higher security while instead a hidden sense of anxiety, of restlessness, and of emptiness is diffused ever more in spite of the systematic use on a large scale and in all the domains of the anesthetic spirituality recently invented. From the sense of ``being'', of stability, also that of the limit, as principle, it could not originate in a natural way, even in a more external domain, for the reaffirmation of forces and processes which have become stronger than those that put them recklessly into a movement of temporality.

But to consider the civilization in its entirety, it remains absolutely problematic where solid points of leverage can be found in a civilization like the modern one that is in every way, to an extent unprecedented in the past, a civilization of time. On the other hand, it appears quite obvious that, more than a correction, there would be the end of a form, the origination of a new form. Thus, logically, at best, we can think only of different orientations in some particular domain and especially in that which, almost like an awakening, a few differentiated men can still propose and implement invisibly.\footnote{The differentiated type, defined by the possession of the dimension of being, is contrasted to the existential orientations suitable to a period of dissolution, like the current one, as we formulated in \emph{Ride the Tiger}.}

\flrightit{Posted on 2019-08-25 by Aeneas}

